%%%%%%%%%%%%%%%%%%%%%%%%%%%%%%%%%%%%%%%%%%%%%%%%%%%%%%%%%%%%%%%%%%%%%%%%%%%%%%
\chapter{Manifestly Covariant Hamiltonian Dynamics}
%%%%%%%%%%%%%%%%%%%%%%%%%%%%%%%%%%%%%%%%%%%%%%%%%%%%%%%%%%%%%%%%%%%%%%%%%%%%%%

\chapterprecis{
This chapter presents a systematic review of the manifestly covariant
Hamiltonian dynamics developed by Stueckelberg and subsequently
generalized by Horwitz and Piron. Particular emphasis is placed upon the
physical motivation of the formalism, its canonical structure, and the
features that make it an appropriate deterministic foundation for a
relativistic stochastic mechanics. The chapter concludes with a critical
comparison with conventional constrained relativistic dynamics.
}

%%%%%%%%%%%%%%%%%%%%%%%%%%%%%%%%%%%%%%%%%%%%%%%%%%%%%%%%%%%%%%%%%%%%%%%%%%%%%%
\section{Chapter Overview}
%%%%%%%%%%%%%%%%%%%%%%%%%%%%%%%%%%%%%%%%%%%%%%%%%%%%%%%%%%%%%%%%%%%%%%%%%%%%%%

Manifestly covariant Hamiltonian dynamics occupies a unique position
among formulations of relativistic mechanics. Unlike conventional
approaches based upon mass-shell constraints, the theory describes the
motion of particles through ordinary Hamilton equations generated by an
invariant Hamiltonian function and an independent evolution parameter.

From a mathematical point of view, the resulting formalism is
remarkably close to ordinary Hamiltonian mechanics. Canonical variables,
Poisson brackets, Hamilton equations, canonical transformations,
Hamilton–Jacobi theory, Liouville dynamics and canonical quantization
retain essentially the same structure as in nonrelativistic mechanics.

The physical interpretation, however, differs profoundly. All four
space-time coordinates become dynamical variables, while evolution is
generated with respect to an invariant parameter rather than the
coordinate time.

These features make the SHP formalism particularly attractive as the
deterministic basis of a relativistic stochastic mechanics.

The objectives of the present chapter are fourfold.

First, we review the historical development of the theory and the
physical problems that motivated its introduction.

Second, we derive the Hamiltonian formulation in a systematic manner,
emphasizing its relation to conventional relativistic dynamics.

Third, we examine the principal consequences of the formalism,
including off-shell dynamics, many-particle systems, and interactions
with external electromagnetic fields.

Finally, we assess the extent to which the SHP theory satisfies the
requirements formulated in Chapter~1 and identify the issues that remain
open.




%%%%%%%%%%%%%%%%%%%%%%%%%%%%%%%%%%%%%%%%%%%%%%%%%%%%%%%%%%%%%%%%%%%%%%%%%%%%%%
\section{Physical Motivation}
%%%%%%%%%%%%%%%%%%%%%%%%%%%%%%%%%%%%%%%%%%%%%%%%%%%%%%%%%%%%%%%%%%%%%%%%%%%%%%

Hamiltonian mechanics occupies a central position in theoretical
physics. Its geometric structure underlies classical mechanics,
statistical mechanics, quantum mechanics and quantum field theory.
Whenever possible, new physical theories are formulated within the
Hamiltonian framework, thereby preserving canonical variables, Poisson
brackets, conservation laws and variational principles.

It is therefore natural to ask whether relativistic particle dynamics
can also be formulated as an unconstrained Hamiltonian system.

At first sight the answer appears obvious. Since special relativity is
well understood, one might expect that Hamiltonian mechanics can be made
relativistic simply by replacing three-dimensional vectors with
four-vectors.

The situation is considerably more subtle.

The origin of the difficulty lies in the fundamentally different role
played by time in Newtonian and relativistic mechanics.

%%%%%%%%%%%%%%%%%%%%%%%%%%%%%%%%%%%%%%%%%%%%%%%%%%%%%%%%%%%%%%%%%%%%%%%%%%%%%%
\subsection{Hamiltonian Mechanics in Newtonian Physics}
%%%%%%%%%%%%%%%%%%%%%%%%%%%%%%%%%%%%%%%%%%%%%%%%%%%%%%%%%%%%%%%%%%%%%%%%%%%%%%

In nonrelativistic mechanics the configuration of a particle is
specified by the spatial coordinates

\[
q^i(t),
\qquad
i=1,2,3,
\]

while the evolution parameter

\[
t
\]

is an external variable.

The canonical momenta

\[
p_i
=
\frac{\partial L}
{\partial \dot q^i}
\]

together with the Hamiltonian

\[
H(q,p,t)
\]

determine the evolution through Hamilton's equations,

\[
\dot q^i
=
\frac{\partial H}{\partial p_i},
\]

\[
\dot p_i
=
-
\frac{\partial H}{\partial q^i}.
\]

Several important features should be emphasized.

First, time is not a dynamical variable.

Second, all particles evolve with respect to the same universal time.

Third, the Hamiltonian generates physical evolution.

These three properties are responsible for the extraordinary success of
Hamiltonian mechanics.

%%%%%%%%%%%%%%%%%%%%%%%%%%%%%%%%%%%%%%%%%%%%%%%%%%%%%%%%%%%%%%%%%%%%%%%%%%%%%%
\subsection{Attempting a Relativistic Generalization}
%%%%%%%%%%%%%%%%%%%%%%%%%%%%%%%%%%%%%%%%%%%%%%%%%%%%%%%%%%%%%%%%%%%%%%%%%%%%%%

Special relativity combines space and time into the four-vector

\[
x^\mu
=
(ct,\mathbf{x}),
\]

suggesting that all four coordinates should be treated equally.

However, the conventional Hamiltonian formalism continues to regard the
time coordinate as an external parameter.

This asymmetry immediately creates conceptual difficulties.

If

\[
x^0
\]

is promoted to an ordinary coordinate, then one must also introduce its
canonically conjugate momentum,

\[
p_0.
\]

But if

\[
x^0
\]

is itself a dynamical variable, with respect to what parameter does the
system evolve?

Conversely, if coordinate time remains the evolution parameter, then
space and time are no longer treated symmetrically despite the Lorentz
structure of Minkowski space.

Thus the straightforward relativistic generalization of Hamiltonian
mechanics encounters a fundamental obstacle.

%%%%%%%%%%%%%%%%%%%%%%%%%%%%%%%%%%%%%%%%%%%%%%%%%%%%%%%%%%%%%%%%%%%%%%%%%%%%%%
\subsection{Two Distinct Roles of Time}
%%%%%%%%%%%%%%%%%%%%%%%%%%%%%%%%%%%%%%%%%%%%%%%%%%%%%%%%%%%%%%%%%%%%%%%%%%%%%%

The difficulty may be understood more clearly by recognizing that the
word "time" refers to two conceptually different quantities.

The first is the temporal coordinate,

\[
x^0,
\]

which specifies the position of an event in space-time.

The second is the evolution parameter that orders successive states of
the dynamical system.

In Newtonian mechanics these two concepts coincide.

The coordinate time simultaneously labels events and generates
evolution.

In relativity they need not be identical.

Indeed, once all four coordinates become dynamical variables, the
evolution parameter must necessarily be distinct from the temporal
coordinate.

Failure to distinguish these two concepts is responsible for many of the
conceptual difficulties encountered in relativistic dynamics.

%%%%%%%%%%%%%%%%%%%%%%%%%%%%%%%%%%%%%%%%%%%%%%%%%%%%%%%%%%%%%%%%%%%%%%%%%%%%%%
\subsection{Requirements for a Covariant Hamiltonian Theory}
%%%%%%%%%%%%%%%%%%%%%%%%%%%%%%%%%%%%%%%%%%%%%%%%%%%%%%%%%%%%%%%%%%%%%%%%%%%%%%

The preceding discussion suggests that a satisfactory relativistic
Hamiltonian theory should satisfy the following requirements.

\begin{enumerate}

\item
All four space-time coordinates should be treated as dynamical
variables.

\item
The evolution parameter should be independent of the space-time
coordinates.

\item
Hamilton's equations should retain their canonical form.

\item
The theory should remain manifestly Lorentz covariant.

\item
The nonrelativistic limit should reproduce ordinary Hamiltonian
mechanics.

\end{enumerate}

Remarkably, these requirements almost determine the structure of the
theory uniquely.

%%%%%%%%%%%%%%%%%%%%%%%%%%%%%%%%%%%%%%%%%%%%%%%%%%%%%%%%%%%%%%%%%%%%%%%%%%%%%%
\subsection{The Stueckelberg Idea}
%%%%%%%%%%%%%%%%%%%%%%%%%%%%%%%%%%%%%%%%%%%%%%%%%%%%%%%%%%%%%%%%%%%%%%%%%%%%%%

The decisive step was taken by Stueckelberg, who proposed separating the
two roles traditionally played by time.

Instead of identifying the temporal coordinate with the evolution
parameter, he introduced an independent invariant parameter

\[
\tau,
\]

while promoting

\[
x^0
\]

to an ordinary canonical coordinate.

The particle trajectory is therefore written as

\[
x^\mu=x^\mu(\tau),
\]

and the dynamics is generated by a Hamiltonian

\[
K(x,p)
\]

through ordinary canonical equations.

The resulting theory preserves the geometric symmetry of Minkowski
space while retaining the canonical structure of Hamiltonian mechanics.

This simple conceptual change forms the foundation of the entire SHP
formalism.

%%%%%%%%%%%%%%%%%%%%%%%%%%%%%%%%%%%%%%%%%%%%%%%%%%%%%%%%%%%%%%%%%%%%%%%%%%%%%%
\subsection{Outlook}
%%%%%%%%%%%%%%%%%%%%%%%%%%%%%%%%%%%%%%%%%%%%%%%%%%%%%%%%%%%%%%%%%%%%%%%%%%%%%%

Having identified the physical motivation, we now proceed to construct
the formalism explicitly.

The next section develops the variational principle and derives the
Hamiltonian equations of motion from first principles.



%%%%%%%%%%%%%%%%%%%%%%%%%%%%%%%%%%%%%%%%%%%%%%%%%%%%%%%%%%%%%%%%%%%%%%%%%%%%%%
\section{Variational Principle and Canonical Equations}
%%%%%%%%%%%%%%%%%%%%%%%%%%%%%%%%%%%%%%%%%%%%%%%%%%%%%%%%%%%%%%%%%%%%%%%%%%%%%%

The Hamiltonian formulation of classical mechanics is most naturally
developed from a variational principle. The same philosophy applies to
manifestly covariant dynamics. Rather than postulating the equations of
motion, we seek an action whose stationary points determine the particle
trajectories.

The fundamental dynamical variables are the canonical coordinates

\[
x^\mu(\tau),
\]

and their conjugate momenta

\[
p_\mu(\tau),
\]

where the invariant parameter \(\tau\) labels successive dynamical
states of the system.

%%%%%%%%%%%%%%%%%%%%%%%%%%%%%%%%%%%%%%%%%%%%%%%%%%%%%%%%%%%%%%%%%%%%%%%%%%%%%%
\subsection{Canonical Action}
%%%%%%%%%%%%%%%%%%%%%%%%%%%%%%%%%%%%%%%%%%%%%%%%%%%%%%%%%%%%%%%%%%%%%%%%%%%%%%

The action is chosen in the canonical first-order form

\[
\boxed{
S[x,p]
=
\int_{\tau_1}^{\tau_2}
\left(
p_\mu \dot x^\mu
-
K(x,p)
\right)
d\tau ,
}
\]

where

\[
\dot x^\mu
\equiv
\frac{dx^\mu}{d\tau},
\]

and

\[
K(x,p)
\]

is a Lorentz-scalar Hamiltonian.

The first term defines the canonical symplectic structure, while the
second determines the dynamics.

The action is invariant under Lorentz transformations provided that
\(K\) is a scalar function of the canonical variables.

%%%%%%%%%%%%%%%%%%%%%%%%%%%%%%%%%%%%%%%%%%%%%%%%%%%%%%%%%%%%%%%%%%%%%%%%%%%%%%
\subsection{Variation with Respect to the Coordinates}
%%%%%%%%%%%%%%%%%%%%%%%%%%%%%%%%%%%%%%%%%%%%%%%%%%%%%%%%%%%%%%%%%%%%%%%%%%%%%%

Consider an arbitrary variation

\[
x^\mu
\rightarrow
x^\mu+\delta x^\mu,
\]

subject to

\[
\delta x^\mu(\tau_1)
=
\delta x^\mu(\tau_2)
=
0.
\]

The corresponding variation of the action is

\[
\delta_x S
=
\int
\left(
p_\mu
\delta\dot x^\mu
-
\frac{\partial K}{\partial x^\mu}
\delta x^\mu
\right)
d\tau .
\]

Integrating the first term by parts,

\[
\int
p_\mu
\delta\dot x^\mu
d\tau
=
\left.
p_\mu
\delta x^\mu
\right|_{\tau_1}^{\tau_2}
-
\int
\dot p_\mu
\delta x^\mu
d\tau .
\]

The boundary contribution vanishes because the end points are fixed.

Therefore,

\[
\delta_x S
=
-
\int
\left(
\dot p_\mu
+
\frac{\partial K}{\partial x^\mu}
\right)
\delta x^\mu
d\tau .
\]

Since the variations are arbitrary,

\[
\boxed{
\dot p_\mu
=
-
\frac{\partial K}{\partial x^\mu}.
}
\]

This is the second Hamilton equation.

%%%%%%%%%%%%%%%%%%%%%%%%%%%%%%%%%%%%%%%%%%%%%%%%%%%%%%%%%%%%%%%%%%%%%%%%%%%%%%
\subsection{Variation with Respect to the Momenta}
%%%%%%%%%%%%%%%%%%%%%%%%%%%%%%%%%%%%%%%%%%%%%%%%%%%%%%%%%%%%%%%%%%%%%%%%%%%%%%

Next consider independent variations

\[
p_\mu
\rightarrow
p_\mu+\delta p_\mu .
\]

The action changes according to

\[
\delta_p S
=
\int
\left(
\dot x^\mu
-
\frac{\partial K}{\partial p_\mu}
\right)
\delta p_\mu
d\tau .
\]

Again the variations are arbitrary.

Therefore,

\[
\boxed{
\dot x^\mu
=
\frac{\partial K}{\partial p_\mu}.
}
\]

This is the first Hamilton equation.

%%%%%%%%%%%%%%%%%%%%%%%%%%%%%%%%%%%%%%%%%%%%%%%%%%%%%%%%%%%%%%%%%%%%%%%%%%%%%%
\subsection{Canonical Hamilton Equations}
%%%%%%%%%%%%%%%%%%%%%%%%%%%%%%%%%%%%%%%%%%%%%%%%%%%%%%%%%%%%%%%%%%%%%%%%%%%%%%

Collecting the previous results, we obtain

\[
\boxed{
\frac{dx^\mu}{d\tau}
=
\frac{\partial K}{\partial p_\mu},
}
\]

\[
\boxed{
\frac{dp_\mu}{d\tau}
=
-
\frac{\partial K}{\partial x^\mu}.
}
\]

These equations have exactly the same mathematical structure as the
ordinary Hamilton equations of classical mechanics.

The essential difference is that the evolution parameter is the
Lorentz-invariant scalar \(\tau\), while all four space-time coordinates
are treated as canonical variables.

%%%%%%%%%%%%%%%%%%%%%%%%%%%%%%%%%%%%%%%%%%%%%%%%%%%%%%%%%%%%%%%%%%%%%%%%%%%%%%
\subsection{Conservation of the Hamiltonian}
%%%%%%%%%%%%%%%%%%%%%%%%%%%%%%%%%%%%%%%%%%%%%%%%%%%%%%%%%%%%%%%%%%%%%%%%%%%%%%

Suppose the Hamiltonian possesses no explicit dependence upon the
evolution parameter,

\[
\frac{\partial K}{\partial\tau}=0.
\]

Differentiating \(K\) along the trajectory,

\[
\frac{dK}{d\tau}
=
\frac{\partial K}{\partial x^\mu}
\dot x^\mu
+
\frac{\partial K}{\partial p_\mu}
\dot p_\mu .
\]

Substituting Hamilton's equations,

\[
\frac{dK}{d\tau}
=
\frac{\partial K}{\partial x^\mu}
\frac{\partial K}{\partial p_\mu}
-
\frac{\partial K}{\partial p_\mu}
\frac{\partial K}{\partial x^\mu}
=
0.
\]

Hence,

\[
\boxed{
\frac{dK}{d\tau}=0,
}
\]

so that the Hamiltonian is conserved along every dynamical trajectory.

This conservation law is the direct analogue of energy conservation in
ordinary Hamiltonian mechanics.

%%%%%%%%%%%%%%%%%%%%%%%%%%%%%%%%%%%%%%%%%%%%%%%%%%%%%%%%%%%%%%%%%%%%%%%%%%%%%%
\subsection{Discussion}
%%%%%%%%%%%%%%%%%%%%%%%%%%%%%%%%%%%%%%%%%%%%%%%%%%%%%%%%%%%%%%%%%%%%%%%%%%%%%%

Several important observations should be made.

First, no mass-shell constraint has been imposed in deriving the
equations of motion. The canonical structure follows solely from the
variational principle.

Second, the temporal coordinate \(x^0\) enters on exactly the same
footing as the spatial coordinates. The distinction between space and
time appears only through the Minkowski metric entering the Hamiltonian.

Finally, the formal similarity with ordinary Hamiltonian mechanics is
complete. Once the invariant evolution parameter has been introduced,
all of the familiar machinery of canonical dynamics follows without
modification.



%%%%%%%%%%%%%%%%%%%%%%%%%%%%%%%%%%%%%%%%%%%%%%%%%%%%%%%%%%%%%%%%%%%%%%%%%%%%%%
\section{Poisson Brackets and Canonical Structure}
%%%%%%%%%%%%%%%%%%%%%%%%%%%%%%%%%%%%%%%%%%%%%%%%%%%%%%%%%%%%%%%%%%%%%%%%%%%%%%

The Hamilton equations derived in the previous section determine the
evolution of the canonical variables. It is convenient, however, to
express the dynamics in a form that applies to an arbitrary dynamical
quantity. This naturally leads to the introduction of the Poisson
bracket.

%%%%%%%%%%%%%%%%%%%%%%%%%%%%%%%%%%%%%%%%%%%%%%%%%%%%%%%%%%%%%%%%%%%%%%%%%%%%%%
\subsection{Evolution of an Arbitrary Observable}
%%%%%%%%%%%%%%%%%%%%%%%%%%%%%%%%%%%%%%%%%%%%%%%%%%%%%%%%%%%%%%%%%%%%%%%%%%%%%%

Let

\[
F(x,p,\tau)
\]

be an arbitrary differentiable function on the phase space.

Its total derivative along a dynamical trajectory is

\[
\frac{dF}{d\tau}
=
\frac{\partial F}{\partial x^\mu}
\frac{dx^\mu}{d\tau}
+
\frac{\partial F}{\partial p_\mu}
\frac{dp_\mu}{d\tau}
+
\frac{\partial F}{\partial\tau}.
\]

Substituting Hamilton's equations,

\[
\frac{dx^\mu}{d\tau}
=
\frac{\partial K}{\partial p_\mu},
\]

\[
\frac{dp_\mu}{d\tau}
=
-
\frac{\partial K}{\partial x^\mu},
\]

one obtains

\[
\frac{dF}{d\tau}
=
\frac{\partial F}{\partial x^\mu}
\frac{\partial K}{\partial p_\mu}
-
\frac{\partial F}{\partial p_\mu}
\frac{\partial K}{\partial x^\mu}
+
\frac{\partial F}{\partial\tau}.
\]

This expression suggests introducing a bilinear operation between two
functions on phase space.

%%%%%%%%%%%%%%%%%%%%%%%%%%%%%%%%%%%%%%%%%%%%%%%%%%%%%%%%%%%%%%%%%%%%%%%%%%%%%%
\subsection{Definition of the Poisson Bracket}
%%%%%%%%%%%%%%%%%%%%%%%%%%%%%%%%%%%%%%%%%%%%%%%%%%%%%%%%%%%%%%%%%%%%%%%%%%%%%%

The Poisson bracket of two functions

\[
F(x,p),
\qquad
G(x,p),
\]

is defined by

\[
\boxed{
\{F,G\}
=
\frac{\partial F}{\partial x^\mu}
\frac{\partial G}{\partial p_\mu}
-
\frac{\partial F}{\partial p_\mu}
\frac{\partial G}{\partial x^\mu}.
}
\]

The summation extends over all four space-time components,

\[
\mu=0,1,2,3.
\]

Apart from the extension to four-dimensional phase space, this
definition is identical to the Poisson bracket of ordinary Hamiltonian
mechanics.

%%%%%%%%%%%%%%%%%%%%%%%%%%%%%%%%%%%%%%%%%%%%%%%%%%%%%%%%%%%%%%%%%%%%%%%%%%%%%%
\subsection{Hamilton's Equations}
%%%%%%%%%%%%%%%%%%%%%%%%%%%%%%%%%%%%%%%%%%%%%%%%%%%%%%%%%%%%%%%%%%%%%%%%%%%%%%

The canonical equations may now be written in the compact form

\[
\boxed{
\frac{dx^\mu}{d\tau}
=
\{x^\mu,K\},
}
\]

and

\[
\boxed{
\frac{dp_\mu}{d\tau}
=
\{p_\mu,K\}.
}
\]

More generally,

\[
\boxed{
\frac{dF}{d\tau}
=
\{F,K\}
+
\frac{\partial F}{\partial\tau}.
}
\]

Thus the Hamiltonian generates the evolution of every observable on the
phase space.

%%%%%%%%%%%%%%%%%%%%%%%%%%%%%%%%%%%%%%%%%%%%%%%%%%%%%%%%%%%%%%%%%%%%%%%%%%%%%%
\subsection{Fundamental Canonical Relations}
%%%%%%%%%%%%%%%%%%%%%%%%%%%%%%%%%%%%%%%%%%%%%%%%%%%%%%%%%%%%%%%%%%%%%%%%%%%%%%

Applying the definition to the canonical variables immediately yields

\[
\boxed{
\{x^\mu,x^\nu\}=0,
}
\]

\[
\boxed{
\{p_\mu,p_\nu\}=0,
}
\]

and

\[
\boxed{
\{x^\mu,p_\nu\}
=
\delta^\mu_{\;\nu}.
}
\]

These relations constitute the canonical algebra of the SHP theory.

They are formally identical to the canonical relations of ordinary
Hamiltonian mechanics.

%%%%%%%%%%%%%%%%%%%%%%%%%%%%%%%%%%%%%%%%%%%%%%%%%%%%%%%%%%%%%%%%%%%%%%%%%%%%%%
\subsection{Algebraic Properties}
%%%%%%%%%%%%%%%%%%%%%%%%%%%%%%%%%%%%%%%%%%%%%%%%%%%%%%%%%%%%%%%%%%%%%%%%%%%%%%

The Poisson bracket possesses the standard algebraic properties.

For arbitrary functions

\[
F,\;G,\;H,
\]

one has

\[
\{F,G\}
=
-
\{G,F\},
\]

(antisymmetry),

\[
\{F,G+H\}
=
\{F,G\}
+
\{F,H\},
\]

(linearity),

and

\[
\{F,GH\}
=
\{F,G\}H
+
G\{F,H\},
\]

(the Leibniz rule).

Finally, the brackets satisfy the Jacobi identity,

\[
\boxed{
\{F,\{G,H\}\}
+
\{G,\{H,F\}\}
+
\{H,\{F,G\}\}
=
0.
}
\]

The Jacobi identity guarantees the internal consistency of Hamiltonian
evolution and plays a fundamental role in canonical quantization.

%%%%%%%%%%%%%%%%%%%%%%%%%%%%%%%%%%%%%%%%%%%%%%%%%%%%%%%%%%%%%%%%%%%%%%%%%%%%%%
\subsection{Constants of Motion}
%%%%%%%%%%%%%%%%%%%%%%%%%%%%%%%%%%%%%%%%%%%%%%%%%%%%%%%%%%%%%%%%%%%%%%%%%%%%%%

A quantity

\[
F(x,p)
\]

is conserved whenever

\[
\frac{dF}{d\tau}=0.
\]

If

\[
\frac{\partial F}{\partial\tau}=0,
\]

this condition becomes

\[
\boxed{
\{F,K\}=0.
}
\]

Thus every conserved quantity Poisson-commutes with the Hamiltonian.

Conversely, every observable whose Poisson bracket with the Hamiltonian
vanishes remains constant along the trajectory.

This criterion provides a powerful method for identifying conserved
quantities without explicitly solving the equations of motion.

%%%%%%%%%%%%%%%%%%%%%%%%%%%%%%%%%%%%%%%%%%%%%%%%%%%%%%%%%%%%%%%%%%%%%%%%%%%%%%
\subsection{Physical Interpretation}
%%%%%%%%%%%%%%%%%%%%%%%%%%%%%%%%%%%%%%%%%%%%%%%%%%%%%%%%%%%%%%%%%%%%%%%%%%%%%%

The Poisson bracket determines the infinitesimal canonical flow
generated by the Hamiltonian. Every observable evolves according to the
same geometric law,

\[
\frac{dF}{d\tau}
=
\{F,K\},
\]

exactly as in ordinary Hamiltonian mechanics.

Consequently, the introduction of the invariant evolution parameter does
not alter the canonical structure of the theory. It merely enlarges the
configuration space from three-dimensional space to four-dimensional
space-time.

This remarkable structural stability is one of the principal reasons
that the SHP formalism provides such a natural foundation for a
relativistic stochastic mechanics.



%%%%%%%%%%%%%%%%%%%%%%%%%%%%%%%%%%%%%%%%%%%%%%%%%%%%%%%%%%%%%%%%%%%%%%%%%%%%%%
\section{Canonical Transformations}
%%%%%%%%%%%%%%%%%%%%%%%%%%%%%%%%%%%%%%%%%%%%%%%%%%%%%%%%%%%%%%%%%%%%%%%%%%%%%%

One of the principal advantages of Hamiltonian mechanics is that its
equations preserve their form under a large class of transformations of
the canonical variables. These canonical transformations play a central
role in perturbation theory, Hamilton--Jacobi theory and canonical
quantization.

The same structure exists in the SHP formalism.

%%%%%%%%%%%%%%%%%%%%%%%%%%%%%%%%%%%%%%%%%%%%%%%%%%%%%%%%%%%%%%%%%%%%%%%%%%%%%%
\subsection{Infinitesimal Canonical Transformations}
%%%%%%%%%%%%%%%%%%%%%%%%%%%%%%%%%%%%%%%%%%%%%%%%%%%%%%%%%%%%%%%%%%%%%%%%%%%%%%

Consider an infinitesimal transformation generated by a scalar function

\[
G(x,p,\tau).
\]

The canonical variables change according to

\[
\boxed{
\delta x^\mu
=
\varepsilon
\{x^\mu,G\},
}
\]

\[
\boxed{
\delta p_\mu
=
\varepsilon
\{p_\mu,G\},
}
\]

where

\[
\varepsilon
\]

is an infinitesimal parameter.

Using the fundamental Poisson brackets,

\[
\{x^\mu,p_\nu\}
=
\delta^\mu_{\;\nu},
\]

one immediately finds

\[
\boxed{
\delta x^\mu
=
\varepsilon
\frac{\partial G}{\partial p_\mu},
}
\]

and

\[
\boxed{
\delta p_\mu
=
-
\varepsilon
\frac{\partial G}{\partial x^\mu}.
}
\]

Thus every differentiable scalar function on phase space generates a
canonical transformation.

%%%%%%%%%%%%%%%%%%%%%%%%%%%%%%%%%%%%%%%%%%%%%%%%%%%%%%%%%%%%%%%%%%%%%%%%%%%%%%
\subsection{Preservation of Hamilton's Equations}
%%%%%%%%%%%%%%%%%%%%%%%%%%%%%%%%%%%%%%%%%%%%%%%%%%%%%%%%%%%%%%%%%%%%%%%%%%%%%%

The defining property of a canonical transformation is the invariance of
Hamilton's equations.

To first order in the infinitesimal parameter,

\[
x^\mu
\rightarrow
x^\mu+\delta x^\mu,
\]

\[
p_\mu
\rightarrow
p_\mu+\delta p_\mu,
\]

while the Hamiltonian transforms as

\[
K
\rightarrow
K+\delta K.
\]

Using the Jacobi identity,

\[
\{F,\{G,K\}\}
+
\{G,\{K,F\}\}
+
\{K,\{F,G\}\}
=
0,
\]

one verifies directly that

\[
\frac{dF'}{d\tau}
=
\{F',K'\}
+
\frac{\partial F'}{\partial\tau},
\]

has exactly the same form as before the transformation.

Hamiltonian evolution is therefore invariant under canonical changes of
variables.

%%%%%%%%%%%%%%%%%%%%%%%%%%%%%%%%%%%%%%%%%%%%%%%%%%%%%%%%%%%%%%%%%%%%%%%%%%%%%%
\subsection{Conservation Laws and Symmetries}
%%%%%%%%%%%%%%%%%%%%%%%%%%%%%%%%%%%%%%%%%%%%%%%%%%%%%%%%%%%%%%%%%%%%%%%%%%%%%%

Canonical transformations provide the natural language for describing
continuous symmetries.

Suppose the Hamiltonian satisfies

\[
\boxed{
\{G,K\}=0.
}
\]

Then

\[
\frac{dG}{d\tau}=0,
\]

and the generator

\[
G
\]

is conserved.

Conversely, every conserved quantity generates an infinitesimal symmetry
of the Hamiltonian flow.

This result is the Hamiltonian form of Noether's theorem.

%%%%%%%%%%%%%%%%%%%%%%%%%%%%%%%%%%%%%%%%%%%%%%%%%%%%%%%%%%%%%%%%%%%%%%%%%%%%%%
\subsection{Space-Time Translations}
%%%%%%%%%%%%%%%%%%%%%%%%%%%%%%%%%%%%%%%%%%%%%%%%%%%%%%%%%%%%%%%%%%%%%%%%%%%%%%

Consider the generator

\[
G=a^\mu p_\mu,
\]

where

\[
a^\mu
\]

is a constant four-vector.

The induced transformation is

\[
\delta x^\mu
=
a^\mu,
\]

\[
\delta p_\mu
=
0.
\]

Thus the canonical momentum generates translations in space-time.

If the Hamiltonian possesses translational symmetry,

\[
\frac{\partial K}{\partial x^\mu}=0,
\]

then

\[
p_\mu
\]

is conserved.

%%%%%%%%%%%%%%%%%%%%%%%%%%%%%%%%%%%%%%%%%%%%%%%%%%%%%%%%%%%%%%%%%%%%%%%%%%%%%%
\subsection{Lorentz Transformations}
%%%%%%%%%%%%%%%%%%%%%%%%%%%%%%%%%%%%%%%%%%%%%%%%%%%%%%%%%%%%%%%%%%%%%%%%%%%%%%

Lorentz transformations are generated by

\[
\boxed{
J^{\mu\nu}
=
x^\mu p^\nu
-
x^\nu p^\mu.
}
\]

The infinitesimal transformation is

\[
\delta x^\mu
=
\frac12
\omega_{\alpha\beta}
\{x^\mu,J^{\alpha\beta}\},
\]

where

\[
\omega_{\alpha\beta}
=
-
\omega_{\beta\alpha}
\]

are the infinitesimal rotation parameters.

Using the canonical brackets,

\[
\boxed{
\delta x^\mu
=
\omega^\mu_{\;\nu}
x^\nu,
}
\]

and similarly,

\[
\boxed{
\delta p^\mu
=
\omega^\mu_{\;\nu}
p^\nu.
}
\]

Therefore the canonical formalism is manifestly Lorentz covariant.

%%%%%%%%%%%%%%%%%%%%%%%%%%%%%%%%%%%%%%%%%%%%%%%%%%%%%%%%%%%%%%%%%%%%%%%%%%%%%%
\subsection{Poincaré Algebra}
%%%%%%%%%%%%%%%%%%%%%%%%%%%%%%%%%%%%%%%%%%%%%%%%%%%%%%%%%%%%%%%%%%%%%%%%%%%%%%

The generators

\[
P^\mu=p^\mu,
\]

and

\[
J^{\mu\nu}
=
x^\mu p^\nu
-
x^\nu p^\mu
\]

satisfy the Poincaré algebra under the Poisson bracket.

In particular,

\[
\boxed{
\{P^\mu,P^\nu\}=0,
}
\]

\[
\boxed{
\{J^{\mu\nu},P^\lambda\}
=
\eta^{\mu\lambda}P^\nu
-
\eta^{\nu\lambda}P^\mu,
}
\]

and

\[
\boxed{
\begin{aligned}
\{J^{\mu\nu},J^{\rho\sigma}\}
&=
\eta^{\mu\rho}J^{\nu\sigma}
-
\eta^{\mu\sigma}J^{\nu\rho}
\\
&
-
\eta^{\nu\rho}J^{\mu\sigma}
+
\eta^{\nu\sigma}J^{\mu\rho}.
\end{aligned}
}
\]

These relations demonstrate that the SHP phase space provides a
Hamiltonian realization of the Poincaré group.

%%%%%%%%%%%%%%%%%%%%%%%%%%%%%%%%%%%%%%%%%%%%%%%%%%%%%%%%%%%%%%%%%%%%%%%%%%%%%%
\subsection{Discussion}
%%%%%%%%%%%%%%%%%%%%%%%%%%%%%%%%%%%%%%%%%%%%%%%%%%%%%%%%%%%%%%%%%%%%%%%%%%%%%%

Canonical transformations demonstrate that the SHP formalism is not
merely a covariant rewriting of relativistic mechanics. It possesses the
complete canonical structure of Hamiltonian dynamics, including
symmetries, conservation laws, and the Poincaré algebra.

Consequently, essentially every construction familiar from classical
Hamiltonian mechanics admits a manifestly covariant analogue within the
SHP framework.


%%%%%%%%%%%%%%%%%%%%%%%%%%%%%%%%%%%%%%%%%%%%%%%%%%%%%%%%%%%%%%%%%%%%%%%%%%%%%%
\section{Phase Space Geometry and Symplectic Structure}
%%%%%%%%%%%%%%%%%%%%%%%%%%%%%%%%%%%%%%%%%%%%%%%%%%%%%%%%%%%%%%%%%%%%%%%%%%%%%%

The Hamiltonian formulation developed in the preceding sections possesses
a geometric structure that is independent of the particular Hamiltonian.
This structure is determined entirely by the canonical variables and
their Poisson brackets. Understanding this geometry is essential for the
development of statistical mechanics and stochastic dynamics in later
chapters.

Although the language of symplectic geometry is often regarded as
abstract, its physical content is straightforward. It provides the
natural geometric description of Hamiltonian evolution.

%%%%%%%%%%%%%%%%%%%%%%%%%%%%%%%%%%%%%%%%%%%%%%%%%%%%%%%%%%%%%%%%%%%%%%%%%%%%%%
\subsection{The Covariant Phase Space}
%%%%%%%%%%%%%%%%%%%%%%%%%%%%%%%%%%%%%%%%%%%%%%%%%%%%%%%%%%%%%%%%%%%%%%%%%%%%%%

For a single particle the configuration space is the four-dimensional
Minkowski space

\[
M.
\]

The corresponding phase space is the cotangent bundle

\[
T^{*}M,
\]

whose points are specified by the canonical coordinates

\[
(x^\mu,p_\mu).
\]

Since there are four coordinates and four conjugate momenta, the phase
space is eight-dimensional.

Unlike configuration space, phase space possesses no distinguished
metric structure. Instead, its geometry is determined by the canonical
symplectic form introduced below.

%%%%%%%%%%%%%%%%%%%%%%%%%%%%%%%%%%%%%%%%%%%%%%%%%%%%%%%%%%%%%%%%%%%%%%%%%%%%%%
\subsection{The Canonical One-Form}
%%%%%%%%%%%%%%%%%%%%%%%%%%%%%%%%%%%%%%%%%%%%%%%%%%%%%%%%%%%%%%%%%%%%%%%%%%%%%%

The canonical action

\[
S
=
\int
(p_\mu \dot x^\mu-K)\,d\tau
\]

contains the differential one-form

\[
\boxed{
\theta
=
p_\mu\,dx^\mu.
}
\]

This object is known as the canonical one-form or the Poincaré one-form.

It depends only upon the canonical variables and is independent of the
Hamiltonian.

The canonical action may therefore be written compactly as

\[
S
=
\int
(\theta-K\,d\tau).
\]

Thus the entire dynamical content of Hamiltonian mechanics is encoded in
the canonical one-form together with the Hamiltonian function.

%%%%%%%%%%%%%%%%%%%%%%%%%%%%%%%%%%%%%%%%%%%%%%%%%%%%%%%%%%%%%%%%%%%%%%%%%%%%%%
\subsection{The Symplectic Two-Form}
%%%%%%%%%%%%%%%%%%%%%%%%%%%%%%%%%%%%%%%%%%%%%%%%%%%%%%%%%%%%%%%%%%%%%%%%%%%%%%

Taking the exterior derivative of the canonical one-form gives

\[
\boxed{
\omega
=
d\theta
=
dp_\mu\wedge dx^\mu .
}
\]

This closed two-form is called the symplectic form.

Several important observations follow immediately.

First,

\[
d\omega
=
0,
\]

since

\[
d^{\,2}=0.
\]

Second,

\[
\omega
\]

is independent of the Hamiltonian.

Consequently, every Hamiltonian system constructed on the same phase
space possesses the same underlying symplectic geometry.

%%%%%%%%%%%%%%%%%%%%%%%%%%%%%%%%%%%%%%%%%%%%%%%%%%%%%%%%%%%%%%%%%%%%%%%%%%%%%%
\subsection{Hamiltonian Vector Field}
%%%%%%%%%%%%%%%%%%%%%%%%%%%%%%%%%%%%%%%%%%%%%%%%%%%%%%%%%%%%%%%%%%%%%%%%%%%%%%

The Hamiltonian determines a vector field on phase space.

Using Hamilton's equations,

\[
\dot x^\mu
=
\frac{\partial K}{\partial p_\mu},
\]

\[
\dot p_\mu
=
-
\frac{\partial K}{\partial x^\mu},
\]

one introduces the Hamiltonian vector field

\[
\boxed{
X_K
=
\frac{\partial K}{\partial p_\mu}
\frac{\partial}{\partial x^\mu}
-
\frac{\partial K}{\partial x^\mu}
\frac{\partial}{\partial p_\mu}.
}
\]

The integral curves of this vector field are precisely the dynamical
trajectories of the particle.

Different Hamiltonians correspond to different vector fields, but all of
them are defined on the same symplectic manifold.

%%%%%%%%%%%%%%%%%%%%%%%%%%%%%%%%%%%%%%%%%%%%%%%%%%%%%%%%%%%%%%%%%%%%%%%%%%%%%%
\subsection{Hamiltonian Flow}
%%%%%%%%%%%%%%%%%%%%%%%%%%%%%%%%%%%%%%%%%%%%%%%%%%%%%%%%%%%%%%%%%%%%%%%%%%%%%%

The solution of Hamilton's equations generates a one-parameter family of
transformations

\[
\Phi_\tau :
T^{*}M
\rightarrow
T^{*}M,
\]

called the Hamiltonian flow.

If the initial point in phase space is

\[
(x^\mu_0,p_{\mu0}),
\]

then

\[
\Phi_\tau
(x^\mu_0,p_{\mu0})
=
(x^\mu(\tau),p_\mu(\tau)).
\]

Thus Hamiltonian dynamics should be viewed as a continuous flow on phase
space rather than merely as a collection of differential equations.

%%%%%%%%%%%%%%%%%%%%%%%%%%%%%%%%%%%%%%%%%%%%%%%%%%%%%%%%%%%%%%%%%%%%%%%%%%%%%%
\subsection{Preservation of the Symplectic Structure}
%%%%%%%%%%%%%%%%%%%%%%%%%%%%%%%%%%%%%%%%%%%%%%%%%%%%%%%%%%%%%%%%%%%%%%%%%%%%%%

One of the fundamental properties of Hamiltonian evolution is that the
symplectic form is preserved along every trajectory.

Symbolically,

\[
\boxed{
\Phi_\tau^{*}\omega
=
\omega,
}
\]

where

\[
\Phi_\tau^{*}
\]

denotes the pull-back generated by the Hamiltonian flow.

This result is the geometric expression of the fact that Hamiltonian
evolution is a canonical transformation.

Unlike an arbitrary coordinate transformation, Hamiltonian evolution
does not distort the intrinsic geometry of phase space.

%%%%%%%%%%%%%%%%%%%%%%%%%%%%%%%%%%%%%%%%%%%%%%%%%%%%%%%%%%%%%%%%%%%%%%%%%%%%%%
\subsection{Physical Interpretation}
%%%%%%%%%%%%%%%%%%%%%%%%%%%%%%%%%%%%%%%%%%%%%%%%%%%%%%%%%%%%%%%%%%%%%%%%%%%%%%

The preservation of the symplectic form has a direct physical meaning.

Infinitesimal phase-space cells may change their shape during the
evolution, but their symplectic structure remains unchanged.

Consequently, Hamiltonian motion is reversible and conserves the
fundamental geometric structure of phase space.

This property ultimately leads to Liouville's theorem and the existence
of an invariant phase-space measure.

%%%%%%%%%%%%%%%%%%%%%%%%%%%%%%%%%%%%%%%%%%%%%%%%%%%%%%%%%%%%%%%%%%%%%%%%%%%%%%
\subsection{Connection with Statistical Mechanics}
%%%%%%%%%%%%%%%%%%%%%%%%%%%%%%%%%%%%%%%%%%%%%%%%%%%%%%%%%%%%%%%%%%%%%%%%%%%%%%

The geometric structure developed above provides the natural setting for
statistical mechanics.

Probability densities are defined on phase space,

\[
\rho(x,p,\tau),
\]

rather than on configuration space alone.

Their evolution is generated by the Hamiltonian flow.

The preservation of the symplectic structure guarantees the existence of
a conserved phase-space volume element, which forms the basis of the
Liouville equation derived in the next section.

Later in this monograph, the same geometric framework will be used to
construct relativistic stochastic dynamics. In that case the
deterministic Hamiltonian flow generated by

\[
X_K
\]

will be supplemented by stochastic fluctuations, leading to a diffusion
process on the covariant phase space.



%%%%%%%%%%%%%%%%%%%%%%%%%%%%%%%%%%%%%%%%%%%%%%%%%%%%%%%%%%%%%%%%%%%%%%%%%%%%%%
\section{Liouville's Theorem and Phase-Space Dynamics}
%%%%%%%%%%%%%%%%%%%%%%%%%%%%%%%%%%%%%%%%%%%%%%%%%%%%%%%%%%%%%%%%%%%%%%%%%%%%%%

The Hamiltonian equations determine the evolution of individual
trajectories in phase space. In many physical applications, however,
one is interested not in a single trajectory but in an ensemble of
systems. The corresponding statistical description requires the
evolution of probability densities rather than individual particles.

The key result making such a description possible is Liouville's
theorem, which states that Hamiltonian evolution preserves the natural
volume measure on phase space.

%%%%%%%%%%%%%%%%%%%%%%%%%%%%%%%%%%%%%%%%%%%%%%%%%%%%%%%%%%%%%%%%%%%%%%%%%%%%%%
\subsection{Hamiltonian Flow in Phase Space}
%%%%%%%%%%%%%%%%%%%%%%%%%%%%%%%%%%%%%%%%%%%%%%%%%%%%%%%%%%%%%%%%%%%%%%%%%%%%%%

The state of a particle is represented by a point

\[
z^A=(x^\mu,p_\mu),
\qquad A=1,\ldots,8,
\]

in the eight-dimensional phase space.

Hamilton's equations may be written compactly as

\[
\frac{dz^A}{d\tau}
=
V^A(z),
\]

where

\[
V^A
=
\left(
\frac{\partial K}{\partial p_\mu},
-
\frac{\partial K}{\partial x^\mu}
\right)
\]

is the Hamiltonian vector field.

The evolution of the system is therefore represented by a continuous
flow generated by the vector field \(V^A\).

%%%%%%%%%%%%%%%%%%%%%%%%%%%%%%%%%%%%%%%%%%%%%%%%%%%%%%%%%%%%%%%%%%%%%%%%%%%%%%
\subsection{Divergence of the Hamiltonian Flow}
%%%%%%%%%%%%%%%%%%%%%%%%%%%%%%%%%%%%%%%%%%%%%%%%%%%%%%%%%%%%%%%%%%%%%%%%%%%%%%

The divergence of the phase-space velocity field is

\[
\partial_A V^A
=
\frac{\partial}{\partial x^\mu}
\left(
\frac{\partial K}{\partial p_\mu}
\right)
+
\frac{\partial}{\partial p_\mu}
\left(
-
\frac{\partial K}{\partial x^\mu}
\right).
\]

Since mixed partial derivatives commute,

\[
\frac{\partial^2K}
{\partial x^\mu\partial p_\mu}
=
\frac{\partial^2K}
{\partial p_\mu\partial x^\mu},
\]

the two contributions cancel exactly.

Hence

\[
\boxed{
\partial_A V^A
=
0.
}
\]

The Hamiltonian flow is therefore divergence-free.

This simple identity is the analytical content of Liouville's theorem.

%%%%%%%%%%%%%%%%%%%%%%%%%%%%%%%%%%%%%%%%%%%%%%%%%%%%%%%%%%%%%%%%%%%%%%%%%%%%%%
\subsection{Conservation of Phase-Space Volume}
%%%%%%%%%%%%%%%%%%%%%%%%%%%%%%%%%%%%%%%%%%%%%%%%%%%%%%%%%%%%%%%%%%%%%%%%%%%%%%

Consider an infinitesimal volume element

\[
d\Gamma
=
d^4x\,d^4p.
\]

Because the divergence of the Hamiltonian vector field vanishes,

\[
\boxed{
\frac{d}{d\tau}(d\Gamma)=0.
}
\]

Thus the phase-space volume occupied by an ensemble of systems remains
constant during the evolution.

Individual volume elements may deform, stretch or rotate, but their
total volume is preserved.

This property is completely analogous to the incompressible flow of an
ideal fluid.

%%%%%%%%%%%%%%%%%%%%%%%%%%%%%%%%%%%%%%%%%%%%%%%%%%%%%%%%%%%%%%%%%%%%%%%%%%%%%%
\subsection{Liouville Equation}
%%%%%%%%%%%%%%%%%%%%%%%%%%%%%%%%%%%%%%%%%%%%%%%%%%%%%%%%%%%%%%%%%%%%%%%%%%%%%%

Let

\[
\rho(x,p,\tau)
\]

denote the probability density on phase space.

Probability conservation requires

\[
\frac{\partial\rho}{\partial\tau}
+
\partial_A
(\rho V^A)
=
0.
\]

Using the divergence-free property,

\[
\partial_A V^A=0,
\]

this becomes

\[
\boxed{
\frac{\partial\rho}{\partial\tau}
+
V^A\partial_A\rho
=
0.
}
\]

Substituting the explicit form of the Hamiltonian vector field gives

\[
\boxed{
\frac{\partial\rho}{\partial\tau}
+
\frac{\partial K}{\partial p_\mu}
\frac{\partial\rho}{\partial x^\mu}
-
\frac{\partial K}{\partial x^\mu}
\frac{\partial\rho}{\partial p_\mu}
=
0.
}
\]

Using the Poisson bracket, this equation assumes the compact form

\[
\boxed{
\frac{\partial\rho}{\partial\tau}
+
\{\rho,K\}
=
0.
}
\]

This is the covariant Liouville equation.

%%%%%%%%%%%%%%%%%%%%%%%%%%%%%%%%%%%%%%%%%%%%%%%%%%%%%%%%%%%%%%%%%%%%%%%%%%%%%%
\subsection{Alternative Geometric Derivation}
%%%%%%%%%%%%%%%%%%%%%%%%%%%%%%%%%%%%%%%%%%%%%%%%%%%%%%%%%%%%%%%%%%%%%%%%%%%%%%

The preceding derivation relied upon the divergence of the Hamiltonian
vector field.

The same result follows directly from the symplectic geometry developed
in the previous section.

Hamiltonian evolution preserves the symplectic form,

\[
\Phi_\tau^*\omega=\omega.
\]

Consequently, the associated volume form

\[
\Omega
=
\frac{\omega^4}{4!}
\]

is also preserved,

\[
\Phi_\tau^*\Omega=\Omega.
\]

Since

\[
\Omega
=
d^4x\,d^4p,
\]

the phase-space volume remains invariant under Hamiltonian evolution.

This geometric argument establishes Liouville's theorem without
reference to any particular coordinate system.

%%%%%%%%%%%%%%%%%%%%%%%%%%%%%%%%%%%%%%%%%%%%%%%%%%%%%%%%%%%%%%%%%%%%%%%%%%%%%%
\subsection{Expectation Values}
%%%%%%%%%%%%%%%%%%%%%%%%%%%%%%%%%%%%%%%%%%%%%%%%%%%%%%%%%%%%%%%%%%%%%%%%%%%%%%

For any observable

\[
F(x,p),
\]

the ensemble average is defined by

\[
\boxed{
\langle F\rangle
=
\int
F(x,p)\,
\rho(x,p,\tau)\,
d^4x\,d^4p.
}
\]

Using the Liouville equation and assuming that surface terms vanish at
infinity, one obtains

\[
\frac{d}{d\tau}
\langle F\rangle
=
\langle
\{F,K\}
\rangle.
\]

This relation is the statistical counterpart of Hamilton's equations.

It shows that the deterministic evolution of observables extends
naturally to statistical ensembles.

%%%%%%%%%%%%%%%%%%%%%%%%%%%%%%%%%%%%%%%%%%%%%%%%%%%%%%%%%%%%%%%%%%%%%%%%%%%%%%
\subsection{Physical Interpretation}
%%%%%%%%%%%%%%%%%%%%%%%%%%%%%%%%%%%%%%%%%%%%%%%%%%%%%%%%%%%%%%%%%%%%%%%%%%%%%%

Liouville's theorem establishes that deterministic Hamiltonian dynamics
is an incompressible flow on phase space.

No probability is created or destroyed during the evolution; it is
merely transported by the Hamiltonian flow.

This property explains why Hamiltonian mechanics provides the natural
foundation of statistical mechanics.

Equally importantly for the purposes of the present work, it also
provides the deterministic limit of relativistic stochastic dynamics.

When stochastic fluctuations are introduced, the Liouville equation will
be replaced by a relativistic Fokker--Planck equation. The latter may be
viewed as a deformation of the deterministic transport equation derived
above, with additional diffusion terms describing the effects of random
motion.


%%%%%%%%%%%%%%%%%%%%%%%%%%%%%%%%%%%%%%%%%%%%%%%%%%%%%%%%%%%%%%%%%%%%%%%%%%%%%%
\section{Hamilton--Jacobi Theory}
%%%%%%%%%%%%%%%%%%%%%%%%%%%%%%%%%%%%%%%%%%%%%%%%%%%%%%%%%%%%%%%%%%%%%%%%%%%%%%

The Hamilton--Jacobi formulation provides an alternative representation
of Hamiltonian mechanics in which the dynamics is described by a single
scalar function rather than by Hamilton's equations. Besides its
importance in classical mechanics, this formulation occupies a central
position in the transition to quantum mechanics and will play a crucial
role in the stochastic theory developed later in this monograph.

%%%%%%%%%%%%%%%%%%%%%%%%%%%%%%%%%%%%%%%%%%%%%%%%%%%%%%%%%%%%%%%%%%%%%%%%%%%%%%
\subsection{Hamilton's Principal Function}
%%%%%%%%%%%%%%%%%%%%%%%%%%%%%%%%%%%%%%%%%%%%%%%%%%%%%%%%%%%%%%%%%%%%%%%%%%%%%%

Consider the action evaluated along a physical trajectory,

\[
S=S(x^\mu,\tau),
\]

viewed as a function of the final space-time point and the evolution
parameter.

The differential of the action is

\[
dS
=
p_\mu\,dx^\mu
-
K\,d\tau.
\]

Comparison with the total differential

\[
dS
=
\frac{\partial S}{\partial x^\mu}dx^\mu
+
\frac{\partial S}{\partial\tau}d\tau
\]

immediately yields the fundamental relations

\[
\boxed{
p_\mu
=
\frac{\partial S}{\partial x^\mu},
}
\]

and

\[
\boxed{
\frac{\partial S}{\partial\tau}
=
-
K.
}
\]

These equations identify the momentum as the gradient of the principal
function.

%%%%%%%%%%%%%%%%%%%%%%%%%%%%%%%%%%%%%%%%%%%%%%%%%%%%%%%%%%%%%%%%%%%%%%%%%%%%%%
\subsection{Derivation of the Hamilton--Jacobi Equation}
%%%%%%%%%%%%%%%%%%%%%%%%%%%%%%%%%%%%%%%%%%%%%%%%%%%%%%%%%%%%%%%%%%%%%%%%%%%%%%

The Hamiltonian is a function of the canonical variables,

\[
K=K(x,p).
\]

Replacing the momentum by

\[
p_\mu
=
\frac{\partial S}{\partial x^\mu},
\]

one obtains

\[
K
=
K
\left(
x,
\frac{\partial S}{\partial x}
\right).
\]

Substituting into the previous equation gives the Hamilton--Jacobi
equation

\[
\boxed{
\frac{\partial S}{\partial\tau}
+
K
\left(
x,
\frac{\partial S}{\partial x}
\right)
=
0.
}
\]

This single nonlinear first-order partial differential equation is
completely equivalent to Hamilton's equations.

%%%%%%%%%%%%%%%%%%%%%%%%%%%%%%%%%%%%%%%%%%%%%%%%%%%%%%%%%%%%%%%%%%%%%%%%%%%%%%
\subsection{Free Particle}
%%%%%%%%%%%%%%%%%%%%%%%%%%%%%%%%%%%%%%%%%%%%%%%%%%%%%%%%%%%%%%%%%%%%%%%%%%%%%%

For a free particle,

\[
K
=
\frac{p^\mu p_\mu}{2M},
\]

the Hamilton--Jacobi equation becomes

\[
\boxed{
\frac{\partial S}{\partial\tau}
+
\frac{1}{2M}
\frac{\partial S}{\partial x^\mu}
\frac{\partial S}{\partial x_\mu}
=
0.
}
\]

Using the Minkowski metric

\[
\eta_{\mu\nu}
=
\mathrm{diag}(+,-,-,-),
\]

this may be written explicitly as

\[
\frac{\partial S}{\partial\tau}
+
\frac{1}{2M}
\left[
\left(
\frac{\partial S}{\partial x^0}
\right)^2
-
\left(
\nabla S
\right)^2
\right]
=
0.
\]

This equation is the manifestly covariant analogue of the ordinary
Hamilton--Jacobi equation.

%%%%%%%%%%%%%%%%%%%%%%%%%%%%%%%%%%%%%%%%%%%%%%%%%%%%%%%%%%%%%%%%%%%%%%%%%%%%%%
\subsection{Characteristics}
%%%%%%%%%%%%%%%%%%%%%%%%%%%%%%%%%%%%%%%%%%%%%%%%%%%%%%%%%%%%%%%%%%%%%%%%%%%%%%

The Hamilton--Jacobi equation determines the dynamics through its
characteristic curves.

Differentiating

\[
p_\mu
=
\frac{\partial S}{\partial x^\mu},
\]

with respect to the evolution parameter and using the Hamilton--Jacobi
equation reproduces Hamilton's equations,

\[
\dot x^\mu
=
\frac{\partial K}{\partial p_\mu},
\]

\[
\dot p_\mu
=
-
\frac{\partial K}{\partial x^\mu}.
\]

Thus the Hamilton--Jacobi equation and the canonical equations are
strictly equivalent descriptions of the same dynamics.

%%%%%%%%%%%%%%%%%%%%%%%%%%%%%%%%%%%%%%%%%%%%%%%%%%%%%%%%%%%%%%%%%%%%%%%%%%%%%%
\subsection{Complete Integrals}
%%%%%%%%%%%%%%%%%%%%%%%%%%%%%%%%%%%%%%%%%%%%%%%%%%%%%%%%%%%%%%%%%%%%%%%%%%%%%%

A complete solution of the Hamilton--Jacobi equation depends upon as
many independent constants as there are degrees of freedom.

For the SHP theory this requires four independent integration constants,

\[
\alpha^\mu,
\]

together with an additional additive constant.

Once a complete integral

\[
S(x,\alpha,\tau)
\]

is known, the particle trajectories follow from

\[
\beta_\mu
=
\frac{\partial S}{\partial\alpha^\mu},
\]

where

\[
\beta_\mu
\]

are constants determined by the initial conditions.

The Hamilton--Jacobi equation therefore transforms the integration of
Hamilton's equations into the solution of a single partial differential
equation.

%%%%%%%%%%%%%%%%%%%%%%%%%%%%%%%%%%%%%%%%%%%%%%%%%%%%%%%%%%%%%%%%%%%%%%%%%%%%%%
\subsection{Relation to the Principle of Least Action}
%%%%%%%%%%%%%%%%%%%%%%%%%%%%%%%%%%%%%%%%%%%%%%%%%%%%%%%%%%%%%%%%%%%%%%%%%%%%%%

The Hamilton--Jacobi equation may be regarded as the differential form of
Hamilton's principle.

Instead of determining the trajectory directly by varying the action,
one first determines the principal function

\[
S(x,\tau),
\]

whose level surfaces define the family of classical trajectories.

The entire dynamics is therefore encoded in the geometry of the action
function.

%%%%%%%%%%%%%%%%%%%%%%%%%%%%%%%%%%%%%%%%%%%%%%%%%%%%%%%%%%%%%%%%%%%%%%%%%%%%%%
\subsection{Physical Interpretation}
%%%%%%%%%%%%%%%%%%%%%%%%%%%%%%%%%%%%%%%%%%%%%%%%%%%%%%%%%%%%%%%%%%%%%%%%%%%%%%

The principal function possesses a clear physical interpretation.

Its gradient determines the local four-momentum,

\[
p_\mu
=
\partial_\mu S.
\]

Consequently, the surfaces

\[
S(x,\tau)=\mathrm{constant}
\]

play the role of relativistic wavefronts.

The particle trajectories are orthogonal to these hypersurfaces in the
sense determined by the canonical momentum.

This interpretation provides the bridge between classical mechanics,
geometrical optics, and wave mechanics.

%%%%%%%%%%%%%%%%%%%%%%%%%%%%%%%%%%%%%%%%%%%%%%%%%%%%%%%%%%%%%%%%%%%%%%%%%%%%%%
\subsection{Connection with Quantum Theory}
%%%%%%%%%%%%%%%%%%%%%%%%%%%%%%%%%%%%%%%%%%%%%%%%%%%%%%%%%%%%%%%%%%%%%%%%%%%%%%

The importance of the Hamilton--Jacobi equation extends far beyond
classical mechanics.

If one introduces the eikonal representation

\[
\Psi
=
A
e^{iS/\hbar},
\]

the phase of the wave function is identified with Hamilton's principal
function.

In the limit

\[
\hbar\rightarrow0,
\]

the Klein--Gordon and Dirac equations reduce to the Hamilton--Jacobi
equation for the classical action.

Conversely, the stochastic theory developed later in this monograph will
modify the Hamilton--Jacobi equation by the appearance of additional
diffusion terms.

These corrections ultimately generate the quantum dynamics.

%%%%%%%%%%%%%%%%%%%%%%%%%%%%%%%%%%%%%%%%%%%%%%%%%%%%%%%%%%%%%%%%%%%%%%%%%%%%%%
\subsection{Discussion}
%%%%%%%%%%%%%%%%%%%%%%%%%%%%%%%%%%%%%%%%%%%%%%%%%%%%%%%%%%%%%%%%%%%%%%%%%%%%%%

The Hamilton--Jacobi formulation provides the natural meeting point of
classical mechanics, wave mechanics and stochastic dynamics.

Within the SHP framework it retains exactly the same conceptual role as
in ordinary Hamiltonian mechanics while remaining manifestly Lorentz
covariant.

For this reason it will become one of the principal tools in the
construction of relativistic stochastic mechanics developed in Part~III.



%%%%%%%%%%%%%%%%%%%%%%%%%%%%%%%%%%%%%%%%%%%%%%%%%%%%%%%%%%%%%%%%%%%%%%%%%%%%%%
\section{Free Particle Dynamics}
%%%%%%%%%%%%%%%%%%%%%%%%%%%%%%%%%%%%%%%%%%%%%%%%%%%%%%%%%%%%%%%%%%%%%%%%%%%%%%

The dynamics of a free particle provides the simplest nontrivial
application of the manifestly covariant Hamiltonian formalism.
Although the resulting equations are elementary, they illustrate the
principal conceptual differences between SHP mechanics and conventional
relativistic dynamics.

%%%%%%%%%%%%%%%%%%%%%%%%%%%%%%%%%%%%%%%%%%%%%%%%%%%%%%%%%%%%%%%%%%%%%%%%%%%%%%
\subsection{The Free Hamiltonian}
%%%%%%%%%%%%%%%%%%%%%%%%%%%%%%%%%%%%%%%%%%%%%%%%%%%%%%%%%%%%%%%%%%%%%%%%%%%%%%

In the absence of external interactions the Hamiltonian must satisfy
three fundamental requirements.

First, it must be a Lorentz scalar.

Second, it must depend only upon the canonical momenta.

Third, it should reduce to the ordinary nonrelativistic Hamiltonian in
the appropriate limit.

The unique quadratic Hamiltonian satisfying these conditions is

\[
\boxed{
K_0
=
\frac{p^\mu p_\mu}{2M},
}
\]

where \(M\) is a constant parameter with the dimensions of mass.

Unlike the invariant mass

\[
m^2
=
p^\mu p_\mu,
\]

the parameter \(M\) is fixed once the Hamiltonian has been chosen.

Its physical interpretation will be discussed below.

%%%%%%%%%%%%%%%%%%%%%%%%%%%%%%%%%%%%%%%%%%%%%%%%%%%%%%%%%%%%%%%%%%%%%%%%%%%%%%
\subsection{Hamilton's Equations}
%%%%%%%%%%%%%%%%%%%%%%%%%%%%%%%%%%%%%%%%%%%%%%%%%%%%%%%%%%%%%%%%%%%%%%%%%%%%%%

Since the free Hamiltonian is independent of the coordinates,

\[
\frac{\partial K_0}{\partial x^\mu}=0,
\]

Hamilton's equations immediately give

\[
\boxed{
\frac{dp_\mu}{d\tau}=0.
}
\]

Thus the four-momentum remains constant along the trajectory.

The coordinate equations become

\[
\boxed{
\frac{dx^\mu}{d\tau}
=
\frac{p^\mu}{M}.
}
\]

The four-velocity with respect to the invariant parameter is therefore

\[
u^\mu
=
\frac{dx^\mu}{d\tau}
=
\frac{p^\mu}{M}.
\]

%%%%%%%%%%%%%%%%%%%%%%%%%%%%%%%%%%%%%%%%%%%%%%%%%%%%%%%%%%%%%%%%%%%%%%%%%%%%%%
\subsection{Worldline}
%%%%%%%%%%%%%%%%%%%%%%%%%%%%%%%%%%%%%%%%%%%%%%%%%%%%%%%%%%%%%%%%%%%%%%%%%%%%%%

Since the momentum is constant, the equations integrate immediately,

\[
\boxed{
x^\mu(\tau)
=
x^\mu_0
+
\frac{p^\mu}{M}
(\tau-\tau_0).
}
\]

The free particle therefore moves along a straight worldline in
Minkowski space.

This result is the covariant analogue of uniform rectilinear motion in
Newtonian mechanics.

%%%%%%%%%%%%%%%%%%%%%%%%%%%%%%%%%%%%%%%%%%%%%%%%%%%%%%%%%%%%%%%%%%%%%%%%%%%%%%
\subsection{Relation Between Momentum and Velocity}
%%%%%%%%%%%%%%%%%%%%%%%%%%%%%%%%%%%%%%%%%%%%%%%%%%%%%%%%%%%%%%%%%%%%%%%%%%%%%%

Using

\[
u^\mu
=
\frac{dx^\mu}{d\tau},
\]

one finds

\[
\boxed{
p^\mu
=
Mu^\mu.
}
\]

This relation resembles the familiar Newtonian expression

\[
\mathbf{p}=m\mathbf{v},
\]

but its interpretation is different.

The velocity is defined with respect to the invariant evolution
parameter rather than the coordinate time.

%%%%%%%%%%%%%%%%%%%%%%%%%%%%%%%%%%%%%%%%%%%%%%%%%%%%%%%%%%%%%%%%%%%%%%%%%%%%%%
\subsection{Invariant Mass}
%%%%%%%%%%%%%%%%%%%%%%%%%%%%%%%%%%%%%%%%%%%%%%%%%%%%%%%%%%%%%%%%%%%%%%%%%%%%%%

The invariant mass is defined by

\[
\boxed{
m^2
=
p^\mu p_\mu.
}
\]

Using

\[
p^\mu
=
Mu^\mu,
\]

one obtains

\[
m^2
=
M^2
u^\mu u_\mu.
\]

Since the free momentum is conserved,

\[
\boxed{
\frac{dm}{d\tau}=0.
}
\]

Thus a free particle remains on a single mass shell.

Only interactions can produce off-shell evolution.

%%%%%%%%%%%%%%%%%%%%%%%%%%%%%%%%%%%%%%%%%%%%%%%%%%%%%%%%%%%%%%%%%%%%%%%%%%%%%%
\subsection{The Parameter Mass \(M\)}
%%%%%%%%%%%%%%%%%%%%%%%%%%%%%%%%%%%%%%%%%%%%%%%%%%%%%%%%%%%%%%%%%%%%%%%%%%%%%%

One of the most characteristic features of the SHP theory is the
appearance of the parameter \(M\).

At first sight one might identify \(M\) with the physical mass of the
particle.

Such an interpretation is generally incorrect.

The quantity \(M\) enters the Hamiltonian as a fixed parameter that
determines the scale relating momentum to the derivative with respect to
the invariant evolution parameter.

The physical invariant mass is instead determined dynamically through

\[
m^2
=
p^\mu p_\mu.
\]

For free motion,

\[
m
=
\mathrm{constant},
\]

but interactions may change its value.

Consequently,

\[
M
\]

and

\[
m
\]

play conceptually different roles.

%%%%%%%%%%%%%%%%%%%%%%%%%%%%%%%%%%%%%%%%%%%%%%%%%%%%%%%%%%%%%%%%%%%%%%%%%%%%%%
\subsection{The On-Shell Limit}
%%%%%%%%%%%%%%%%%%%%%%%%%%%%%%%%%%%%%%%%%%%%%%%%%%%%%%%%%%%%%%%%%%%%%%%%%%%%%%

Suppose the initial conditions satisfy

\[
p^\mu p_\mu=M^2.
\]

Since

\[
p^\mu
\]

is conserved,

\[
p^\mu p_\mu
=
M^2
\]

holds for all values of the evolution parameter.

In this case

\[
u^\mu u_\mu
=
1,
\]

and

\[
d\tau
=
ds,
\]

where

\[
ds
\]

denotes the infinitesimal proper time.

Thus conventional relativistic mechanics is recovered as a special case
of the SHP theory.

%%%%%%%%%%%%%%%%%%%%%%%%%%%%%%%%%%%%%%%%%%%%%%%%%%%%%%%%%%%%%%%%%%%%%%%%%%%%%%
\subsection{Beyond the Mass Shell}
%%%%%%%%%%%%%%%%%%%%%%%%%%%%%%%%%%%%%%%%%%%%%%%%%%%%%%%%%%%%%%%%%%%%%%%%%%%%%%

The SHP formalism does not require

\[
p^\mu p_\mu
=
M^2.
\]

Instead,

\[
p^\mu p_\mu
\]

is a dynamical quantity whose value is determined by the equations of
motion.

The free particle remains on whatever mass shell is specified by its
initial conditions, but interacting particles may evolve continuously
between different mass shells.

This possibility is completely absent from conventional constrained
relativistic mechanics.

%%%%%%%%%%%%%%%%%%%%%%%%%%%%%%%%%%%%%%%%%%%%%%%%%%%%%%%%%%%%%%%%%%%%%%%%%%%%%%
\subsection{Discussion}
%%%%%%%%%%%%%%%%%%%%%%%%%%%%%%%%%%%%%%%%%%%%%%%%%%%%%%%%%%%%%%%%%%%%%%%%%%%%%%

The free particle illustrates the essential conceptual innovation of the
SHP formalism.

Hamiltonian evolution is generated with respect to an invariant
parameter, all four space-time coordinates are treated as canonical
variables, and the invariant mass becomes a dynamical quantity rather
than a fixed constraint.

These features distinguish SHP mechanics fundamentally from the
conventional constrained formulation of relativistic particle dynamics
and provide the flexibility required for the stochastic extensions
considered in later chapters.



%%%%%%%%%%%%%%%%%%%%%%%%%%%%%%%%%%%%%%%%%%%%%%%%%%%%%%%%%%%%%%%%%%%%%%%%%%%%%%
\section{The Physical Interpretation of the Parameter Mass}
%%%%%%%%%%%%%%%%%%%%%%%%%%%%%%%%%%%%%%%%%%%%%%%%%%%%%%%%%%%%%%%%%%%%%%%%%%%%%%

The appearance of two quantities with the dimensions of mass is one of
the most distinctive features of the SHP formalism. In conventional
relativistic mechanics the mass is introduced through the constraint

\[
p^\mu p_\mu = m^2,
\]

and no ambiguity arises.

In manifestly covariant Hamiltonian dynamics, however, two different
quantities naturally appear:

\[
M,
\]

which enters the Hamiltonian,

and

\[
m^2 = p^\mu p_\mu,
\]

which characterizes the instantaneous dynamical state of the particle.

Understanding the distinction between these quantities is essential for
the physical interpretation of the theory.

%%%%%%%%%%%%%%%%%%%%%%%%%%%%%%%%%%%%%%%%%%%%%%%%%%%%%%%%%%%%%%%%%%%%%%%%%%%%%%
\subsection{The Origin of the Parameter \(M\)}
%%%%%%%%%%%%%%%%%%%%%%%%%%%%%%%%%%%%%%%%%%%%%%%%%%%%%%%%%%%%%%%%%%%%%%%%%%%%%%

The Hamiltonian of a free particle is

\[
K
=
\frac{p^\mu p_\mu}{2M}.
\]

The appearance of the denominator is not arbitrary.

The Hamiltonian must have dimensions of energy, while

\[
p^\mu p_\mu
\]

has dimensions of momentum squared.

Consequently, a parameter with the dimensions of mass is required.

At the purely dimensional level,

\[
M
\]

plays exactly the same role that the inertial mass plays in the
nonrelativistic Hamiltonian

\[
H
=
\frac{\mathbf p^2}{2m}.
\]

However, the physical interpretation is subtler because the invariant
mass is no longer imposed as a constraint.

%%%%%%%%%%%%%%%%%%%%%%%%%%%%%%%%%%%%%%%%%%%%%%%%%%%%%%%%%%%%%%%%%%%%%%%%%%%%%%
\subsection{Kinematical Role of \(M\)}
%%%%%%%%%%%%%%%%%%%%%%%%%%%%%%%%%%%%%%%%%%%%%%%%%%%%%%%%%%%%%%%%%%%%%%%%%%%%%%

Hamilton's equations give

\[
\frac{dx^\mu}{d\tau}
=
\frac{p^\mu}{M}.
\]

Therefore

\[
M
\]

determines the proportionality between the canonical momentum and the
velocity with respect to the invariant evolution parameter.

It is therefore more accurate to regard

\[
M
\]

as the parameter governing the response of the trajectory to the
Hamiltonian flow.

Unlike

\[
m,
\]

it does not describe the physical state of the particle.

%%%%%%%%%%%%%%%%%%%%%%%%%%%%%%%%%%%%%%%%%%%%%%%%%%%%%%%%%%%%%%%%%%%%%%%%%%%%%%
\subsection{The Dynamical Mass}
%%%%%%%%%%%%%%%%%%%%%%%%%%%%%%%%%%%%%%%%%%%%%%%%%%%%%%%%%%%%%%%%%%%%%%%%%%%%%%

The quantity

\[
m^2
=
p^\mu p_\mu
\]

is determined entirely by the canonical momentum.

Since the momentum evolves according to Hamilton's equations,

\[
m
\]

is itself a dynamical observable.

For a free particle,

\[
\frac{dm}{d\tau}=0,
\]

because

\[
\frac{dp_\mu}{d\tau}=0.
\]

In the presence of interactions,

\[
m
\]

need not remain constant.

Consequently, the SHP formalism naturally admits off-shell motion.

%%%%%%%%%%%%%%%%%%%%%%%%%%%%%%%%%%%%%%%%%%%%%%%%%%%%%%%%%%%%%%%%%%%%%%%%%%%%%%
\subsection{Why \(M\neq m\)}
%%%%%%%%%%%%%%%%%%%%%%%%%%%%%%%%%%%%%%%%%%%%%%%%%%%%%%%%%%%%%%%%%%%%%%%%%%%%%%

It is tempting to identify

\[
M=m,
\]

thereby recovering the familiar relation

\[
p^\mu
=
m u^\mu.
\]

Such an identification is valid only for free particles satisfying the
initial condition

\[
p^\mu p_\mu=M^2.
\]

The SHP formalism deliberately avoids imposing this condition because it
would eliminate the possibility of off-shell evolution.

The distinction between

\[
M
\]

and

\[
m
\]

is therefore not an additional complication but an essential structural
feature of the theory.

%%%%%%%%%%%%%%%%%%%%%%%%%%%%%%%%%%%%%%%%%%%%%%%%%%%%%%%%%%%%%%%%%%%%%%%%%%%%%%
\subsection{The Nonrelativistic Limit}
%%%%%%%%%%%%%%%%%%%%%%%%%%%%%%%%%%%%%%%%%%%%%%%%%%%%%%%%%%%%%%%%%%%%%%%%%%%%%%

The physical meaning of

\[
M
\]

becomes clearer in the nonrelativistic limit.

Suppose

\[
p^\mu
=
(mc,\mathbf p),
\]

with

\[
|\mathbf p|
\ll
mc.
\]

Then

\[
K
=
\frac{m^2c^2-\mathbf p^2}{2M}.
\]

The first term is an additive constant,

while the second gives

\[
K
=
\mathrm{constant}
-
\frac{\mathbf p^2}{2M}.
\]

After subtraction of the irrelevant constant and taking into account the
sign convention of the metric, one recovers the ordinary kinetic energy

\[
\frac{\mathbf p^2}{2M}.
\]

Thus

\[
M
\]

plays exactly the role of the inertial mass appearing in Newtonian
mechanics.

%%%%%%%%%%%%%%%%%%%%%%%%%%%%%%%%%%%%%%%%%%%%%%%%%%%%%%%%%%%%%%%%%%%%%%%%%%%%%%
\subsection{Interpretation as an Inertial Parameter}
%%%%%%%%%%%%%%%%%%%%%%%%%%%%%%%%%%%%%%%%%%%%%%%%%%%%%%%%%%%%%%%%%%%%%%%%%%%%%%

The preceding discussion suggests that

\[
M
\]

should not be regarded as the physical mass measured in scattering
experiments.

Instead,

\[
M
\]

is better interpreted as an inertial parameter defining the Hamiltonian
flow with respect to the invariant evolution parameter.

The measurable invariant mass is

\[
m,
\]

which depends upon the dynamical state of the system.

%%%%%%%%%%%%%%%%%%%%%%%%%%%%%%%%%%%%%%%%%%%%%%%%%%%%%%%%%%%%%%%%%%%%%%%%%%%%%%
\subsection{Role in Interacting Systems}
%%%%%%%%%%%%%%%%%%%%%%%%%%%%%%%%%%%%%%%%%%%%%%%%%%%%%%%%%%%%%%%%%%%%%%%%%%%%%%

The distinction between

\[
M
\]

and

\[
m
\]

becomes particularly important for interacting particles.

Interactions modify the momentum,

\[
p_\mu(\tau),
\]

and therefore alter

\[
m^2=p^\mu p_\mu.
\]

The parameter

\[
M,
\]

however, remains fixed in the Hamiltonian.

Thus the Hamiltonian structure itself is unaffected even though the
particle evolves continuously between different mass shells.

%%%%%%%%%%%%%%%%%%%%%%%%%%%%%%%%%%%%%%%%%%%%%%%%%%%%%%%%%%%%%%%%%%%%%%%%%%%%%%
\subsection{Implications for Quantization}
%%%%%%%%%%%%%%%%%%%%%%%%%%%%%%%%%%%%%%%%%%%%%%%%%%%%%%%%%%%%%%%%%%%%%%%%%%%%%%

The distinction between

\[
M
\]

and

\[
m
\]

persists after canonical quantization.

The Hamiltonian operator contains the fixed parameter

\[
M,
\]

whereas the mass operator is associated with the spectrum of

\[
\hat p^\mu \hat p_\mu.
\]

Consequently, off-shell quantum states arise naturally within the SHP
framework.

Their physical interpretation remains an active subject of research.

%%%%%%%%%%%%%%%%%%%%%%%%%%%%%%%%%%%%%%%%%%%%%%%%%%%%%%%%%%%%%%%%%%%%%%%%%%%%%%
\subsection{Implications for Stochastic Dynamics}
%%%%%%%%%%%%%%%%%%%%%%%%%%%%%%%%%%%%%%%%%%%%%%%%%%%%%%%%%%%%%%%%%%%%%%%%%%%%%%

The distinction between the parameter mass and the dynamical invariant
mass may become even more significant in a stochastic theory.

Random fluctuations modify the particle momentum and therefore influence

\[
m^2
=
p^\mu p_\mu.
\]

The parameter

\[
M,
\]

however, continues to determine the deterministic Hamiltonian flow.

This separation between kinematics and dynamics provides one of the
principal reasons why the SHP formalism appears particularly well suited
for a stochastic extension.

%%%%%%%%%%%%%%%%%%%%%%%%%%%%%%%%%%%%%%%%%%%%%%%%%%%%%%%%%%%%%%%%%%%%%%%%%%%%%%
\subsection{Discussion}
%%%%%%%%%%%%%%%%%%%%%%%%%%%%%%%%%%%%%%%%%%%%%%%%%%%%%%%%%%%%%%%%%%%%%%%%%%%%%%

The introduction of the parameter

\[
M
\]

is not an arbitrary modification of relativistic mechanics.

It is a direct consequence of formulating the theory as an unconstrained
Hamiltonian system with an independent evolution parameter.

The dynamical invariant mass

\[
m
\]

characterizes the instantaneous physical state of the particle, whereas

\[
M
\]

determines the canonical evolution generated by the Hamiltonian.

Recognizing the different physical roles of these two quantities is
essential for understanding both off-shell dynamics and the stochastic
generalization developed in later chapters.


%%%%%%%%%%%%%%%%%%%%%%%%%%%%%%%%%%%%%%%%%%%%%%%%%%%%%%%%%%%%%%%%%%%%%%%%%%%%%%
\section{Off-Shell Dynamics}
%%%%%%%%%%%%%%%%%%%%%%%%%%%%%%%%%%%%%%%%%%%%%%%%%%%%%%%%%%%%%%%%%%%%%%%%%%%%%%

One of the most distinctive consequences of manifestly covariant
Hamiltonian dynamics is that the invariant mass of a particle is no
longer imposed as a fixed kinematical constraint. Instead, it becomes a
dynamical quantity whose evolution is determined by the equations of
motion.

This property is usually referred to as \emph{off-shell dynamics}.

The concept has no analogue in conventional relativistic particle
mechanics, where the mass-shell condition

\[
p^\mu p_\mu = m^2
\]

is imposed before the equations of motion are derived.

%%%%%%%%%%%%%%%%%%%%%%%%%%%%%%%%%%%%%%%%%%%%%%%%%%%%%%%%%%%%%%%%%%%%%%%%%%%%%%
\subsection{The Mass-Shell Constraint in Conventional Mechanics}
%%%%%%%%%%%%%%%%%%%%%%%%%%%%%%%%%%%%%%%%%%%%%%%%%%%%%%%%%%%%%%%%%%%%%%%%%%%%%%

In the standard formulation of relativistic mechanics the four-momentum
is constrained by

\[
\boxed{
p^\mu p_\mu = m^2,
}
\]

where

\[
m
\]

is interpreted as the physical rest mass.

The constraint reduces the number of independent dynamical variables and
restricts the particle to a single hyperboloid in momentum space.

Consequently, the invariant mass is not determined dynamically.

It is specified once and for all before solving the equations of motion.

%%%%%%%%%%%%%%%%%%%%%%%%%%%%%%%%%%%%%%%%%%%%%%%%%%%%%%%%%%%%%%%%%%%%%%%%%%%%%%
\subsection{Absence of the Constraint in SHP}
%%%%%%%%%%%%%%%%%%%%%%%%%%%%%%%%%%%%%%%%%%%%%%%%%%%%%%%%%%%%%%%%%%%%%%%%%%%%%%

The SHP formalism follows a completely different strategy.

The Hamiltonian

\[
K(x,p)
\]

is defined on the full phase space

\[
(x^\mu,p_\mu),
\]

without imposing any algebraic relation between the canonical momenta.

The equations of motion therefore determine the evolution of

\[
p_\mu(\tau)
\]

directly.

The invariant quantity

\[
m^2
=
p^\mu p_\mu
\]

is then calculated from the solution rather than assumed in advance.

%%%%%%%%%%%%%%%%%%%%%%%%%%%%%%%%%%%%%%%%%%%%%%%%%%%%%%%%%%%%%%%%%%%%%%%%%%%%%%
\subsection{Evolution of the Invariant Mass}
%%%%%%%%%%%%%%%%%%%%%%%%%%%%%%%%%%%%%%%%%%%%%%%%%%%%%%%%%%%%%%%%%%%%%%%%%%%%%%

Differentiating

\[
m^2
=
p^\mu p_\mu,
\]

with respect to the invariant parameter,

\[
\frac{d}{d\tau}
(m^2)
=
2p^\mu
\frac{dp_\mu}{d\tau}.
\]

Using Hamilton's equations,

\[
\frac{dp_\mu}{d\tau}
=
-
\frac{\partial K}{\partial x^\mu},
\]

one finds

\[
\boxed{
\frac{d}{d\tau}
(m^2)
=
-2
p^\mu
\frac{\partial K}{\partial x^\mu}.
}
\]

This equation is exact.

It shows immediately that the invariant mass changes whenever the
Hamiltonian possesses explicit dependence upon the space-time
coordinates.

%%%%%%%%%%%%%%%%%%%%%%%%%%%%%%%%%%%%%%%%%%%%%%%%%%%%%%%%%%%%%%%%%%%%%%%%%%%%%%
\subsection{Free Particle}
%%%%%%%%%%%%%%%%%%%%%%%%%%%%%%%%%%%%%%%%%%%%%%%%%%%%%%%%%%%%%%%%%%%%%%%%%%%%%%

For the free Hamiltonian,

\[
K
=
\frac{p^\mu p_\mu}{2M},
\]

one has

\[
\frac{\partial K}{\partial x^\mu}=0,
\]

and therefore

\[
\boxed{
\frac{dm^2}{d\tau}=0.
}
\]

The particle remains on a single mass shell.

Thus free motion reproduces the conventional picture despite the absence
of any imposed constraint.

%%%%%%%%%%%%%%%%%%%%%%%%%%%%%%%%%%%%%%%%%%%%%%%%%%%%%%%%%%%%%%%%%%%%%%%%%%%%%%
\subsection{Interacting Particle}
%%%%%%%%%%%%%%%%%%%%%%%%%%%%%%%%%%%%%%%%%%%%%%%%%%%%%%%%%%%%%%%%%%%%%%%%%%%%%%

Suppose now that

\[
K
=
\frac{p^\mu p_\mu}{2M}
+
V(x).
\]

Then

\[
\boxed{
\frac{d}{d\tau}(m^2)
=
-2
p^\mu
\partial_\mu V.
}
\]

The invariant mass therefore varies whenever the particle moves through
regions where the scalar potential possesses a nonvanishing
space-time gradient.

Off-shell motion is therefore a natural dynamical consequence of the
interaction rather than an additional postulate.

%%%%%%%%%%%%%%%%%%%%%%%%%%%%%%%%%%%%%%%%%%%%%%%%%%%%%%%%%%%%%%%%%%%%%%%%%%%%%%
\subsection{Energy Exchange}
%%%%%%%%%%%%%%%%%%%%%%%%%%%%%%%%%%%%%%%%%%%%%%%%%%%%%%%%%%%%%%%%%%%%%%%%%%%%%%

The variation of

\[
m^2
\]

should not be interpreted as the creation or destruction of matter.

Instead, it reflects an exchange between the invariant mass of the
particle and the interaction energy represented by the Hamiltonian.

Since

\[
K
\]

is conserved whenever it possesses no explicit dependence upon

\[
\tau,
\]

changes in

\[
p^\mu p_\mu
\]

are compensated by corresponding changes in the interaction term.

The total Hamiltonian remains constant.

%%%%%%%%%%%%%%%%%%%%%%%%%%%%%%%%%%%%%%%%%%%%%%%%%%%%%%%%%%%%%%%%%%%%%%%%%%%%%%
\subsection{Physical Interpretation}
%%%%%%%%%%%%%%%%%%%%%%%%%%%%%%%%%%%%%%%%%%%%%%%%%%%%%%%%%%%%%%%%%%%%%%%%%%%%%%

Off-shell motion should be understood as an extension of the classical
concept of energy conservation.

In Newtonian mechanics the kinetic energy varies while the total energy
remains constant.

In SHP mechanics the invariant mass may vary while the Hamiltonian
remains constant.

The invariant mass therefore becomes another dynamical quantity capable
of exchanging energy with the interaction field.

%%%%%%%%%%%%%%%%%%%%%%%%%%%%%%%%%%%%%%%%%%%%%%%%%%%%%%%%%%%%%%%%%%%%%%%%%%%%%%
\subsection{The On-Shell Limit}
%%%%%%%%%%%%%%%%%%%%%%%%%%%%%%%%%%%%%%%%%%%%%%%%%%%%%%%%%%%%%%%%%%%%%%%%%%%%%%

If the interaction satisfies

\[
p^\mu
\partial_\mu V=0,
\]

then

\[
\frac{dm^2}{d\tau}=0.
\]

The particle remains on a fixed mass shell even in the presence of the
interaction.

Conventional relativistic mechanics therefore appears as a special class
of solutions of the more general SHP dynamics.

%%%%%%%%%%%%%%%%%%%%%%%%%%%%%%%%%%%%%%%%%%%%%%%%%%%%%%%%%%%%%%%%%%%%%%%%%%%%%%
\subsection{Discussion}
%%%%%%%%%%%%%%%%%%%%%%%%%%%%%%%%%%%%%%%%%%%%%%%%%%%%%%%%%%%%%%%%%%%%%%%%%%%%%%

The possibility of off-shell evolution is the defining physical
consequence of the unconstrained Hamiltonian formulation.

The invariant mass is no longer prescribed by a kinematical constraint
but becomes a dynamical observable determined by the interaction.

This enlargement of the space of admissible trajectories is one of the
principal reasons why the SHP theory provides a promising framework for
relativistic stochastic dynamics.


%%%%%%%%%%%%%%%%%%%%%%%%%%%%%%%%%%%%%%%%%%%%%%%%%%%%%%%%%%%%%%%%%%%%%%%%%%%%%%
\chapter{Electromagnetic Interaction}
%%%%%%%%%%%%%%%%%%%%%%%%%%%%%%%%%%%%%%%%%%%%%%%%%%%%%%%%%%%%%%%%%%%%%%%%%%%%%%

%%%%%%%%%%%%%%%%%%%%%%%%%%%%%%%%%%%%%%%%%%%%%%%%%%%%%%%%%%%%%%%%%%%%%%%%%%%%%%
\section{Minimal Coupling}
%%%%%%%%%%%%%%%%%%%%%%%%%%%%%%%%%%%%%%%%%%%%%%%%%%%%%%%%%%%%%%%%%%%%%%%%%%%%%%

The interaction of a charged particle with the electromagnetic field
provides the first physically important application of the SHP
Hamiltonian formalism.

The construction is guided by three fundamental requirements.

\begin{enumerate}

\item Lorentz covariance,

\item gauge invariance,

\item reduction to the ordinary Lorentz force in the appropriate limit.

\end{enumerate}

These conditions uniquely determine the interaction Hamiltonian.

%%%%%%%%%%%%%%%%%%%%%%%%%%%%%%%%%%%%%%%%%%%%%%%%%%%%%%%%%%%%%%%%%%%%%%%%%%%%%%
\subsection{Gauge Potential}
%%%%%%%%%%%%%%%%%%%%%%%%%%%%%%%%%%%%%%%%%%%%%%%%%%%%%%%%%%%%%%%%%%%%%%%%%%%%%%

The electromagnetic field is described by the four-potential

\[
A^\mu(x).
\]

Under a gauge transformation,

\[
\boxed{
A_\mu
\rightarrow
A_\mu+\partial_\mu\Lambda.
}
\]

The observable electromagnetic tensor is

\[
\boxed{
F_{\mu\nu}
=
\partial_\mu A_\nu
-
\partial_\nu A_\mu.
}
\]

Since mixed derivatives commute,

\[
F_{\mu\nu}
\]

is gauge invariant.

%%%%%%%%%%%%%%%%%%%%%%%%%%%%%%%%%%%%%%%%%%%%%%%%%%%%%%%%%%%%%%%%%%%%%%%%%%%%%%
\subsection{Minimal Coupling}
%%%%%%%%%%%%%%%%%%%%%%%%%%%%%%%%%%%%%%%%%%%%%%%%%%%%%%%%%%%%%%%%%%%%%%%%%%%%%%

Gauge invariance requires the replacement

\[
\boxed{
p_\mu
\longrightarrow
\pi_\mu
=
p_\mu
-
eA_\mu.
}
\]

The quantity

\[
\pi_\mu
\]

is the kinetic four-momentum.

The canonical momentum remains

\[
p_\mu,
\]

while the physical momentum carried by the particle is

\[
\pi_\mu.
\]

This distinction is essential.

%%%%%%%%%%%%%%%%%%%%%%%%%%%%%%%%%%%%%%%%%%%%%%%%%%%%%%%%%%%%%%%%%%%%%%%%%%%%%%
\subsection{Hamiltonian}
%%%%%%%%%%%%%%%%%%%%%%%%%%%%%%%%%%%%%%%%%%%%%%%%%%%%%%%%%%%%%%%%%%%%%%%%%%%%%%

The gauge-invariant Hamiltonian becomes

\[
\boxed{
K
=
\frac{1}{2M}
(p_\mu-eA_\mu)
(p^\mu-eA^\mu).
}
\]

More generally,

\[
K
=
\frac{\pi^\mu\pi_\mu}{2M}
+
V(x),
\]

if additional scalar interactions are present.

%%%%%%%%%%%%%%%%%%%%%%%%%%%%%%%%%%%%%%%%%%%%%%%%%%%%%%%%%%%%%%%%%%%%%%%%%%%%%%
\subsection{Hamilton's Equations}
%%%%%%%%%%%%%%%%%%%%%%%%%%%%%%%%%%%%%%%%%%%%%%%%%%%%%%%%%%%%%%%%%%%%%%%%%%%%%%

Differentiating the Hamiltonian gives

\[
\boxed{
\frac{dx^\mu}{d\tau}
=
\frac{\pi^\mu}{M}.
}
\]

Hence the kinetic momentum determines the physical velocity.

For the canonical momentum,

\[
\frac{dp_\mu}{d\tau}
=
-
\frac{\partial K}{\partial x^\mu}.
\]

The derivative contains the explicit coordinate dependence of the gauge
potential,

\[
\partial_\mu A_\nu.
\]

After differentiation,

\[
\boxed{
\frac{dp_\mu}{d\tau}
=
\frac{e}{M}
\pi^\nu
\partial_\mu A_\nu.
}
\]

%%%%%%%%%%%%%%%%%%%%%%%%%%%%%%%%%%%%%%%%%%%%%%%%%%%%%%%%%%%%%%%%%%%%%%%%%%%%%%
\subsection{Equation for the Kinetic Momentum}
%%%%%%%%%%%%%%%%%%%%%%%%%%%%%%%%%%%%%%%%%%%%%%%%%%%%%%%%%%%%%%%%%%%%%%%%%%%%%%

The physical equation is obtained from

\[
\pi_\mu
=
p_\mu-eA_\mu.
\]

Differentiating,

\[
\frac{d\pi_\mu}{d\tau}
=
\frac{dp_\mu}{d\tau}
-
e
\frac{dA_\mu}{d\tau}.
\]

Since

\[
\frac{dA_\mu}{d\tau}
=
\frac{dx^\nu}{d\tau}
\partial_\nu A_\mu,
\]

one obtains

\[
\frac{d\pi_\mu}{d\tau}
=
\frac{e}{M}
\pi^\nu
\partial_\mu A_\nu
-
\frac{e}{M}
\pi^\nu
\partial_\nu A_\mu.
\]

Collecting terms,

\[
\boxed{
\frac{d\pi_\mu}{d\tau}
=
\frac{e}{M}
F_{\mu\nu}
\pi^\nu.
}
\]

This is the covariant Lorentz-force equation in SHP mechanics.

%%%%%%%%%%%%%%%%%%%%%%%%%%%%%%%%%%%%%%%%%%%%%%%%%%%%%%%%%%%%%%%%%%%%%%%%%%%%%%
\subsection{Second-Order Equation}
%%%%%%%%%%%%%%%%%%%%%%%%%%%%%%%%%%%%%%%%%%%%%%%%%%%%%%%%%%%%%%%%%%%%%%%%%%%%%%

Using

\[
\pi^\mu
=
M
\frac{dx^\mu}{d\tau},
\]

one immediately finds

\[
\boxed{
M
\frac{d^2x^\mu}{d\tau^2}
=
e
F^\mu{}_\nu
\frac{dx^\nu}{d\tau}.
}
\]

This equation has precisely the same form as the ordinary Lorentz-force
equation except that differentiation is taken with respect to the
invariant evolution parameter.

%%%%%%%%%%%%%%%%%%%%%%%%%%%%%%%%%%%%%%%%%%%%%%%%%%%%%%%%%%%%%%%%%%%%%%%%%%%%%%
\subsection{Gauge Invariance}
%%%%%%%%%%%%%%%%%%%%%%%%%%%%%%%%%%%%%%%%%%%%%%%%%%%%%%%%%%%%%%%%%%%%%%%%%%%%%%

Under the gauge transformation

\[
A_\mu
\rightarrow
A_\mu+\partial_\mu\Lambda,
\]

the canonical momentum changes according to

\[
p_\mu
\rightarrow
p_\mu
+
e
\partial_\mu\Lambda,
\]

so that

\[
\boxed{
\pi_\mu
=
p_\mu-eA_\mu
}
\]

remains invariant.

Consequently,

\[
K,
\]

Hamilton's equations,

and the Lorentz-force equation are all gauge invariant.

Gauge symmetry is therefore incorporated directly into the canonical
Hamiltonian formalism.



